%% kapa: exam=1 round=26 year=2015 period=1 form=A subject=civil_law subject_ko=민법 seq=37 number=37
%% source: Q-Net 감정평가사 기출문제 (cst00309) artlSeq=5138001 file=extracted/1008/5138001/제26회_감정평가사 1차_기출/제26회_감정평가사 1차_기출/1교시_감평_A형.pdf
%% conversion: kordoc 4.24.51 -> tools/gen1.py (auto); review=proofread-sonnet; flags=none
\begin{question}{37}
지상권에 관한 설명으로 옳지 않은 것은?
\begin{choices}
  \item 수목의 소유를 목적으로 한 지상권의 최단존속기간은 30년이다.
  \item 토지의 지상권자는 타인에게 그 권리를 양도하거나 그 권리의 존속기간 내에서 그 토지를 임대할 수 있다.
  \item 지상권이 소멸한 경우에 지상권설정자가 상당한 가액을 제공하여 그 토지에 현존하는 공작물이나 수목의 매수를 청구한 때에는 지상권자는 정당한 이유없이 이를 거절하지 못한다.
  \item 타인의 토지의 지하 또는 지상의 공간을 상하의 범위를 정하여 사용할 수 있는 권리를 물권으로 취득하는 것도 허용된다.
  \item 지상권자가 2년 이상의 지료를 지급하지 아니한 때에는 지상권설정자는 지상권의 소멸을 청구할 수 있으나, 당사자의 약정으로 그 기간을 단축할 수 있다.
\end{choices}
\end{question}
\answer{5}
